\documentclass[11pt]{article}
\usepackage[a4paper,margin=27mm]{geometry}
\usepackage{amsmath,amssymb,amsthm}
\usepackage[T1]{fontenc}
\usepackage{lmodern}
\usepackage{hyperref}
\usepackage{array,longtable}
\hypersetup{colorlinks=true,urlcolor=blue,citecolor=blue,linkcolor=blue}
\newcommand{\Tr}{\operatorname{Tr}}
\newcommand{\ii}{\mathrm{i}}
\newcommand{\KD}{\operatorname{KD}}
\newtheorem{theorem}{Theorem}
\newtheorem{corollary}[theorem]{Corollary}
\newtheorem{proposition}[theorem]{Proposition}
\title{A KD-positive counterexample to the SLD conjecture\\
in arXiv:2511.01996v2, \S6.3 after (6.23), p.~43}
\author{Haobo Ma\\\small ChronoAI Pte Ltd, Singapore
\and Wenlin Zhang\\\small National University of Singapore\\\small ChronoAI Pte Ltd, Singapore}
\date{3 November 2026}
\begin{document}
\maketitle
\begin{abstract}
We refute the mixed-state conjecture of Spriet, Langrenez, Brummelhuis and
De Bi\`evre in arXiv:2511.01996v2, \S6.3 immediately after (6.23), p.~43: a KD-real
state and generator can have a KD-real symmetric logarithmic derivative
with nonzero quantum Fisher information. Two rank-three projections in
dimension nine give a faithful KD-positive state with quantum Fisher
information $1/6$. The same construction gives an explicit one-parameter
family, including two singular mixed endpoints, and counterexamples in
every square dimension $n^2$ with $n\ge3$. We also prove that the exclusion
does hold for every state in the real span of the two basis-diagonal
spaces, for cyclic-prime Fourier bases, and for all KD-real states when
the transition matrix is real up to basis phases. It also holds for
Haar-almost every pair of bases in each finite dimension.
\end{abstract}

\section{The conjecture and conventions}
The unnumbered conjecture in \cite[\S6.3, after (6.23)]{source} asserts that
the symmetric logarithmic derivative (SLD) of a KD-real state under a KD-real
generator cannot itself be KD-real when the quantum Fisher information is
nonzero. Its journal locator is \S6.3, after (168). We give a counterexample
with strictly positive density matrix and strictly positive KD distribution.
For background on KD representations, see \cite[\S II]{survey}.

Let $A=(a_x)$ and $B=(b_k)$ be orthonormal bases with all overlaps nonzero.
The left KD symbol and distribution are, respectively,
\[
 S_C(x,k)=\frac{\langle a_x,Cb_k\rangle}{\langle a_x,b_k\rangle},
 \qquad q_C(x,k)=|\langle a_x,b_k\rangle|^2S_C(x,k).
\]
A Hermitian matrix is KD-real if every $S_C(x,k)$ is real; a density
matrix is KD-positive if every $q_C(x,k)$ is nonnegative.
For $\rho(\theta)=e^{-\ii\theta X}\rho e^{\ii\theta X}$, our normalization is
\begin{equation}\label{eq:sld}
 \frac{L\rho+\rho L}{2}=\ii[\rho,X],\qquad
 I_{\rm QF}=\Tr(\rho L^2).
\end{equation}
These are the defining SLD and QFI conventions of \cite[Appendix E]{source}.
In a density eigenbasis the defining equation gives
$L_{jk}=2\ii(p_j-p_k)X_{jk}/(p_j+p_k)$ when $p_j+p_k>0$;
the illustrative entry formula preceding (6.29) must include this factor $2$.
All formulas below use \eqref{eq:sld}.

\section{Two projections and the counterexample}
All coordinate and frequency labels below are in $\mathbb F_3$. Put
$\omega=(-1+\ii\sqrt3)/2$, and index $\mathbb C^9$ by $(u,v)\in\mathbb F_3^2$.
Take $A$ to be the coordinate basis and set
\[
 F_{(u,v),(r,s)}=\langle a_{u,v},b_{r,s}\rangle
     =\tfrac13\omega^{ur+vs}.
\]
The identity $\sum_{j=0}^2\omega^{mj}=3\,\mathbf1_{m=0}$ proves $F^*F=I$;
all overlaps have modulus $1/3$.
Define matrices
\begin{equation}\label{eq:projections}
 \begin{split}
 P_{(u,v),(a,b)}&=\tfrac13\mathbf1_{v=b}\omega^{v(u-a)},\\
 Q_{(u,v),(a,b)}&=\tfrac13\mathbf1_{u=a}\omega^{u(b-v)}.
 \end{split}
\end{equation}
Each fixed-$v$ block of $P$, and each fixed-$u$ block of $Q$, is a rank-one
orthogonal projection. Thus $P^2=P=P^*$, $Q^2=Q=Q^*$ and
$\Tr P=\Tr Q=3$. Introduce
\begin{equation}\label{eq:family}
 N_\lambda=9+3\lambda,\quad
 \rho_\lambda=\frac{I+\lambda P}{N_\lambda},\quad
 X=\frac{3Q-I}{2},\quad
 H=\ii[P,Q],\quad
 L_\lambda=\frac{3\lambda}{2+\lambda}H.
\end{equation}

\begin{theorem}\label{thm:main}
For every real $\lambda>-1$, $\rho_\lambda$ is a faithful KD-positive
density matrix, and $X$ and $L_\lambda$ are KD-real Hermitian matrices.
The unique SLD in \eqref{eq:sld} is $L_\lambda$, and
\begin{equation}\label{eq:qfi}
 I_{\rm QF}(\rho_\lambda,X)
    =\frac{2\lambda^2}{(2+\lambda)(3+\lambda)}.
\end{equation}
In particular, $\lambda=1$ gives a counterexample with QFI $1/6$ and
density eigenvalues $1/6$ (multiplicity three) and $1/12$ (multiplicity six).
\end{theorem}

\begin{proof}
The eigenvalues of $\rho_\lambda$ are $(1+\lambda)/N_\lambda$ on
$\operatorname{ran}P$ and $1/N_\lambda$ on $\ker P$, so they are positive
and their sum, with multiplicities, is one.

We compute the symbols explicitly. Summing first in the coordinate on
which each projection acts gives
\begin{equation}\label{eq:symbols}
 S_P((u,v),(r,s))=\mathbf1_{r=v},\qquad
 S_Q((u,v),(r,s))=\mathbf1_{s=-u}.
\end{equation}
For the product, direct substitution in \eqref{eq:projections} yields
\[
 (PQ)_{(u,v),(a,b)}=\tfrac19\omega^{v(u-a)+a(b-v)}.
\]
In $(PQF)_{(u,v),(r,s)}$, summing over $b$ forces $a=-s$.
Dividing by $F_{(u,v),(r,s)}$ then gives
\begin{equation}\label{eq:product}
 S_{PQ}=\tfrac13\omega^z,\qquad S_{QP}=\tfrac13\omega^{-z},
 \qquad z=(u+s)(v-r).
\end{equation}
For $QP$, the corresponding sum over $a$ forces $b=r$.
Consequently,
\[
 S_{\rho_\lambda}=\frac{1+\lambda\mathbf1_{r=v}}{N_\lambda},\quad
 S_X=\frac{3\mathbf1_{s=-u}-1}{2},\quad
 S_H=\frac{\ii}{3}(\omega^z-\omega^{-z})\in\{0,\pm1/\sqrt3\}.
\]
These symbols are real, and $S_{\rho_\lambda}>0$. Since
$q_C=S_C/9$ in these bases, $\rho_\lambda$ is strictly KD-positive.
The three matrices are Hermitian because $P,Q$ are Hermitian and
$\ii[P,Q]$ is Hermitian.

For the SLD equation, $PH+HP=H$ follows by expanding the commutator and
using $P^2=P$. Thus
\[
 \frac{\rho_\lambda L_\lambda+L_\lambda\rho_\lambda}{2}
 =\frac{2+\lambda}{2N_\lambda}L_\lambda
 =\frac{3\lambda}{2N_\lambda}\ii[P,Q]
 =\ii[\rho_\lambda,X].
\]
If $D$ is the difference of two SLDs, diagonalizing the positive matrix
$\rho_\lambda$ gives $(p_j+p_k)D_{jk}=0$ with $p_j,p_k>0$, hence $D=0$.

It remains to evaluate the trace, without numerical diagonalization.
Multiplying the product formula above by $P$ gives
\[
 (PQP)_{(u,v),(a,b)}
 =\frac{\omega^{uv-ab}}{27}\sum_{c=0}^2\omega^{2c(b-v)}
 =\tfrac13P_{(u,v),(a,b)}.
\]
Also $(PQ)_{(u,v),(u,v)}=1/9$, so $\Tr(PQ)=1$ and
$\Tr(PQPQ)=1/3$. Expanding $H^2$ gives
\[
 H^2=PQ+QP-PQP-QPQ,\qquad
 \Tr H^2=\tfrac43,\quad \Tr(PH^2)=\tfrac23.
\]
Therefore
\[
 \Tr(\rho_\lambda L_\lambda^2)
 =\frac{9\lambda^2}{(2+\lambda)^2N_\lambda}
       \left(\tfrac43+\tfrac{2\lambda}{3}\right)
 =\frac{2\lambda^2}{(2+\lambda)(3+\lambda)}.
\]
The value is positive whenever $\lambda\ne0$, proving the assertion.
\end{proof}

\begin{corollary}\label{cor:square}
Let $n\ge2$ and $g=\gcd(2,n)$. In dimension $n^2$, use labels in $\mathbb Z/n\mathbb Z$,
$\omega=e^{2\pi\ii/n}$, and the matrices $F,P,Q$ above with $1/3$ replaced
by $1/n$. Set
\[
 \rho_\lambda=\frac{I+\lambda P}{n^2+n\lambda},\quad
 X=\frac{nQ-I}{n-1},\quad
 L_\lambda=\frac{2n\lambda}{(n-1)(2+\lambda)}\ii[P,Q].
\]
For $\lambda>-1$, these give faithful KD-positive states and KD-real
generators and SLDs, with
\[
 I_{\rm QF}=\frac{4(n-g)\lambda^2}{(n-1)^2(2+\lambda)(n+\lambda)}.
\]
For $n\ge3$ and $\lambda\ne0$ this is positive. For odd $n$, it reduces to
$4\lambda^2/((n-1)(2+\lambda)(n+\lambda))$; for $n=2$ it vanishes.
\end{corollary}
\begin{proof}
Character orthogonality gives the two projection-symbol formulas in
\eqref{eq:symbols} unchanged and the product-symbol formulas
$S_{PQ}=\omega^z/n$, $S_{QP}=\omega^{-z}/n$. Thus $\ii[P,Q]$ is KD-real.
The rank of each projection is $n$, so the displayed density is faithful
and normalized, with positive symbol
$(1+\lambda\mathbf1_{r=v})/(n^2+n\lambda)$.
For all $n$, direct multiplication gives
\[
 (PQP)_{(u,v),(a,b)}
 =\frac{\omega^{uv-ab}}{n^3}\sum_{c=0}^{n-1}\omega^{2c(b-v)}.
\]
Its nonzero entries have $2(b-v)=0$ modulo $n$. In the trace
\[
 \Tr(PQPQ)=\sum_{u,v,a,b}(PQP)_{(u,v),(a,b)}Q_{(a,b),(u,v)},
\]
$Q$ forces $a=u$. For each $(u,v)$ there are exactly $g$ choices of $b$,
each contributing $1/n^3$, because the remaining phase is
$\omega^{2u(v-b)}=1$. Hence $\Tr(PQPQ)=g/n$, while
$\Tr(PQ)=1$ as before. Expanding $H^2$ and using trace cyclicity gives
$\Tr H^2=2(n-g)/n$ and $\Tr(PH^2)=(n-g)/n$.
The identity $PH+HP=H$ proves the SLD equation with the displayed
coefficient; positivity proves uniqueness as before. Substitution of these
two traces into $\Tr(\rho_\lambda L_\lambda^2)$ gives the stated formula.
\end{proof}

\begin{corollary}\label{cor:endpoints}
In dimension nine put, for $0\le t\le1$,
\[
 \sigma_t=\frac{t}{3}P+\frac{1-t}{6}(I-P),\qquad
 M_t=\frac{3(3t-1)}{1+t}H.
\]
Then $\sigma_t$ is KD-positive, $M_t$ is a KD-real SLD for the same $X$,
and $I_{\rm QF}=(3t-1)^2/(1+t)$. The state is faithful exactly when
$0<t<1$. The endpoint states have ranks six and three and QFI $1$ and
$2$, respectively; neither is pure.
\end{corollary}
\begin{proof}
The two spectral coefficients are $t/3$ and $(1-t)/6$; their sum is
$(1+t)/6>0$. Since $PH+HP=H$, the Jordan equation on the two spectral
subspaces gives $M_t$ directly, including the endpoints. Its QFI follows
from $\Tr H^2=4/3$ and $\Tr(PH^2)=2/3$.
The symbol of $\sigma_t$ is $t/3$ when $r=v$ and $(1-t)/6$ otherwise.
These numbers are nonnegative, proving KD positivity. For $t<1$ this
is the family in Theorem~\ref{thm:main} with
$\lambda=(3t-1)/(1-t)$; the direct calculation also covers $t=1$.
\end{proof}

\begin{corollary}\label{cor:measurement}
For every square family in Corollary~\ref{cor:square}, both original-basis
probability distributions are uniform and have zero phase derivatives at
$\theta=0$. Thus both classical Fisher informations vanish. When $n\ge3$
and $\lambda\ne0$, measuring the spectral projections of the KD-real SLD
instead attains the positive QFI.
\end{corollary}
\begin{proof}
The diagonals of $P$ in the two bases are $1/n$, and those of
$[P,Q]$ are zero. In the coordinate basis this follows from
$(PQ)_{xx}=(QP)_{xx}=1/n^2$. In the Fourier basis the diagonal of $PQ$
is $n^{-4}\sum_{u,v}\omega^{(u+s)(v-r)}=1/n^2$, since summing over $v$
forces $u=-s$; $QP$ has the same diagonal. Thus the probabilities are
$1/n^2$ and their derivatives vanish. If $E_j$ is an SLD eigenspace
projection with eigenvalue $\ell_j$, then
$\dot p_j=\Tr(E_j\{\rho,L\}/2)=\ell_jp_j$.
Faithfulness gives $p_j>0$, so the measurement's Fisher information is
$\sum_j\dot p_j^2/p_j=\sum_j\ell_j^2p_j=\Tr(\rho L^2)$.
\end{proof}

\section{Mixed-state conditions that do imply exclusion}
\begin{corollary}\label{cor:diagonal}
For any pair of bases with nonzero overlaps, let $\rho$ be a density
matrix in the real linear span of the Hermitian matrices diagonal in
$A$ or in $B$. If $X$ and an SLD $L$ are KD-real, then
$I_{\rm QF}=0$ and $[\rho,X]=0$.
This includes the states $\rho=f_2(A)+g_2(B)$ of source (6.16), even without
restricting $X$ to that form. It also holds whenever $\rho$ is diagonal in either original
basis, without the extra restrictions of the two mixed-state examples in
\cite[\S6.3, pp.~44--45]{source}.
\end{corollary}
\begin{proof}
Write $F_{xk}=\langle a_x,b_k\rangle$. For KD-real Hermitian $X,L$,
\[
 (XL)_{xx}=\sum_k(XF)_{xk}\overline{(LF)_{xk}}
 =\sum_k|F_{xk}|^2S_X(x,k)S_L(x,k)\in\mathbb R.
\]
Thus $[X,L]$ has zero diagonal in $A$. Interchanging $A,B$ conjugates
the Hermitian symbols and proves the same in $B$. Consequently
$\Tr(\rho[X,L])=0$. Multiplying \eqref{eq:sld} by $L$ and taking the trace gives
$\Tr(\rho L^2)=\ii\Tr(\rho[X,L])=0$.
Positivity implies $L\rho^{1/2}=0$ and hence $L\rho=\rho L=0$;
\eqref{eq:sld} now yields $[\rho,X]=0$. No faithfulness is needed.
\end{proof}

\begin{corollary}\label{cor:real}
If the transition matrix can be made real by rephasing basis vectors,
the mixed-state exclusion holds for every KD-real density matrix:
a KD-real SLD has zero QFI. In particular, no counterexample exists
in dimensions one or two.
\end{corollary}
\begin{proof}
Symbols are unchanged by basis rephasings, so assume $F$ is real.
KD reality and nonzero real overlaps make $CF$ real entrywise, so
$C=(CF)F^{\mathsf T}$ is real. Conversely every real Hermitian $C$ is
KD-real. For such $\rho,X,L$ the Jordan side of \eqref{eq:sld} is real,
and the commutator side is purely imaginary. Both vanish, so
$\Tr(\rho L^2)=\Tr(L\{\rho,L\}/2)=0$.
For a $2\times2$ unitary with nonzero entries, rephase its first row
and first column to be positive real. Column orthogonality forces the
remaining entry to be real. In dimension one all matrices commute.
\end{proof}

\begin{corollary}\label{cor:prime}
For the coordinate and Fourier bases of $\mathbb C^{\mathbb Z/p\mathbb Z}$,
with $p$ prime, the mixed-state exclusion holds for every KD-real state.
The same conclusion holds for Haar-almost every pair of bases in any
fixed finite dimension.
\end{corollary}
\begin{proof}
By \cite[Theorems 3.1 and 4.1]{fourier}, the KD-real Hermitian space is
the real span of the pure KD-positive projectors, which correspond to
subgroup cosets and their character modulations. The only subgroups of
$\mathbb Z/p\mathbb Z$ are $\{0\}$ and the whole group. Their projectors
are exactly those of the two original bases. Thus every KD-real state
has the form required by Corollary~\ref{cor:diagonal}.
For the second assertion, \cite[Theorem 1, (7)]{generic} states that the
KD-real Hermitian space equals this same basis-projector span for
Haar-almost every transition matrix. Apply Corollary~\ref{cor:diagonal} again.
\end{proof}

\begin{proposition}\label{prop:algebra}
Suppose a unital complex star-subalgebra $\mathcal A$ of matrices has
KD-real Hermitian part. For every faithful density matrix
$\rho\in\mathcal A$ and every Hermitian $X\in\mathcal A$, the unique SLD
belongs to $\mathcal A$ and is KD-real. Its QFI is positive exactly when
$[\rho,X]\ne0$.
\end{proposition}
\begin{proof}
The map $J_\rho(C)=\{\rho,C\}/2$ preserves the finite-dimensional real
space $\mathcal A_{\rm sa}$. If $\rho\ge mI>0$, then
$\Tr(CJ_\rho(C))=\Tr(\rho C^2)\ge m\Tr(C^2)>0$ for $C\ne0$.
It is therefore invertible on that space, and its inverse sends
$\ii[\rho,X]\in\mathcal A_{\rm sa}$ to the desired $L$.
The same inequality on all Hermitian matrices gives uniqueness and
$I_{\rm QF}>0$ exactly when $L\ne0$, equivalently $[\rho,X]\ne0$.
\end{proof}

The square construction also illustrates this mechanism. With
$x,r\in(\mathbb Z/n\mathbb Z)^2$, put $J(r_1,r_2)=(-r_2,r_1)$ and
$W_ra_x=\omega^{(Jr)\cdot x}a_{x+r}$. Direct application to $a_x$ gives
$W_rW_s=\omega^{(Jr)\cdot s}W_{r+s}$ and $W_r^*=W_{-r}$.
Their span is a star-algebra, and
$S_{W_r}(x,k)=\omega^{(Jr)\cdot x-k\cdot r}$.
In a Hermitian linear combination the coefficients at $r$ and $-r$
are conjugate, so the symbol is real. Indeed the $W_r$ are linearly
independent, since $\Tr(W_r^*W_s)=n^2\mathbf1_{r=s}$.
The averages $n^{-1}\sum_aW_{(a,0)}$ and
$n^{-1}\sum_bW_{(0,b)}$ are precisely $P$ and $Q$.
For $n\ge3$, the generators $W_{(1,0)}$ and $W_{(0,1)}$ do not commute,
as their multiplication phases differ by $\omega^2\ne1$.
Fourier KD-real spaces and their relation to KD-positive states are studied
in \cite{fourier}; here the Jordan-map argument supplies the SLD conclusion.

\section{Coverage and verification}
At $\lambda=1$, $q_{\rho_1}$ takes the values $1/54$ and $1/108$.
Corollary~\ref{cor:measurement} gives the zero classical Fisher information
in the original two bases, consistently with the source Proposition 6.3.

The pure-state conclusion of \cite[Proposition 6.5]{source} retains its
scope. Here $(\rho_1)_{(0,0),(1,0)}=1/36\ne0$, so the basis-diagonal
hypotheses of the two source examples do not apply. The following table
distinguishes the source statements and the extensions established here.
Only \texttt{slbd2026} is a problem identifier for these statements.
\begin{small}
\setlength{\tabcolsep}{4pt}
\begin{longtable}{@{}p{.36\linewidth}p{.23\linewidth}p{.35\linewidth}@{}}
\textbf{Source item} & \textbf{Coverage} & \textbf{Reason}\\\hline
\S6.3 after (6.23), p.~43; all-mixed question, p.~45
 & Refuted-and-corrected & Theorem~\ref{thm:main}; Corollaries~\ref{cor:square}, \ref{cor:endpoints}.\\
Proposition 6.3; phase-insensitive examples (6.16)
 & Prior proof cited & Source proof; Corollary~\ref{cor:measurement} proves the square family's marginals.\\
Lemma 6.4: zero QFI iff commutation
 & Prior proof cited & Source proof; positive-QFI conditions are explicit above.\\
Proposition 6.5: pure-state exclusion
 & Prior proof cited & Unchanged hypotheses; all nontrivial states above are mixed.\\
Two mixed examples, pp.~44--45, after (6.29)
 & Proved & Corollary~\ref{cor:diagonal} removes both examples' extra restrictions.\\
Real transition matrices, pp.~43,45
 & Proved & Corollary~\ref{cor:real} permits every KD-real state.\\
Possibility of a KD-real observable with positive FI, p.~45
 & Proved & Corollary~\ref{cor:measurement}: the KD-real SLD attains QFI.\\
Postselected metrology remark, p.~45, refs.~[42,70]
 & Not covered & Different model and KD distribution; cf.\ \cite{postselected}.\\
Fourier KD-real structure in source ref.~[23]
 & Prior proof cited & \cite{fourier}, \S4; Corollary~\ref{cor:prime} and Proposition~\ref{prop:algebra} give SLD consequences.\\
Generic KD-real structure in source ref.~[69]
 & Prior proof cited & \cite{generic}, Theorem 1; Corollary~\ref{cor:prime} gives Haar-almost-sure exclusion.\\
LCA-group conjectures in source ref.~[24], equations (11)--(12)
 & Not covered & \cite{lca}, p.~7; finite matrix SLD identities do not settle infinite-group closure questions.\\
\end{longtable}
\end{small}
The least counterexample dimension is between three and nine, by
Corollary~\ref{cor:real} and Theorem~\ref{thm:main}; its exact value and a
full classification of bases satisfying mixed-state exclusion are not
determined here.

The pinned Lean file
\href{https://github.com/the-omega-institute/trureturing/blob/71b4592f31a0ce564acedcbee195fbd50f30bd8c/D5/S3/Quantum/Information/KirkwoodDiracSLDRefutation.lean}{\texttt{KirkwoodDiracSLDRefutation.lean}}
contains theorem \texttt{result} of type $\neg\,\mathtt{claim}$ in namespace
\texttt{SprietKirkwoodDiracSLD}, within the module namespace.
It formalizes the refutation with the concrete $\lambda=1$ matrices.
Its \texttt{fisherSLD} convention is one quarter of the QFI in
\eqref{eq:sld}. The linked
\href{https://github.com/the-omega-institute/trureturing/blob/71b4592f31a0ce564acedcbee195fbd50f30bd8c/D5/S3/Quantum/Information/FisherDiscord/Defs.lean}{\texttt{FisherDiscord/Defs.lean}}
defines \texttt{IsSLD} and \texttt{fisherSLD};
\href{https://github.com/the-omega-institute/trureturing/blob/71b4592f31a0ce564acedcbee195fbd50f30bd8c/D5/S3/Quantum/Information/SLDSpectralEnergy.lean}{\texttt{SLDSpectralEnergy.lean}}
defines \texttt{IsDensity} as positive semidefinite with trace one.
The parameter, square-dimension, endpoint, exclusion and algebra results
are written proofs above, not additional declarations of these files.
The accompanying \texttt{verify.py} checks matrix identities over
$\mathbb Q(\sqrt3,\ii)$: six dimension-nine rational parameters,
the two singular endpoints, square orders $n=2,4,6$, and the trace
obstruction in Corollary~\ref{cor:diagonal}. These finite checks are
distinct from the full parameter and arbitrary-dimension proofs.

\begingroup\small
\begin{thebibliography}{9}
\bibitem{source}
M. Spriet, C. Langrenez, R. Brummelhuis and S. De Bi\`evre,
\emph{What Is Special About the Kirkwood--Dirac Distributions?
Only They Produce Natural Conditional Expectations},
Symmetry \textbf{18} (2026), 1008,
\href{https://doi.org/10.3390/sym18061008}{doi:10.3390/sym18061008}.
Posing version: \href{https://arxiv.org/abs/2511.01996v2}{arXiv:2511.01996v2},
\S6.3 after (6.23), p.~43; mixed-state question, p.~45.
\bibitem{survey}
D. R. M. Arvidsson-Shukur, W. F. Braasch Jr., S. De Bi\`evre, J. Dressel,
A. N. Jordan, C. Langrenez, M. Lostaglio, J. S. Lundeen and N. Yunger Halpern,
\emph{Properties and applications of the Kirkwood--Dirac distribution},
New Journal of Physics \textbf{26} (2024), 121201,
\href{https://doi.org/10.1088/1367-2630/ada05d}{doi:10.1088/1367-2630/ada05d};
\href{https://arxiv.org/abs/2403.18899v2}{arXiv:2403.18899v2}, \S II.
\bibitem{fourier}
S. De Bi\`evre, C. Langrenez and D. Radchenko,
\emph{The Kirkwood--Dirac Representation Associated to the Fourier
Transform for Finite Abelian Groups: Positivity},
Annales Henri Poincar\'e \textbf{27} (2026), 3093--3112 (online 2 September 2025),
\href{https://doi.org/10.1007/s00023-025-01614-7}{doi:10.1007/s00023-025-01614-7};
\href{https://arxiv.org/abs/2501.12252v2}{arXiv:2501.12252v2}, Theorems 3.1 and 4.1.
\bibitem{postselected}
S. Das, S. Modak and M. N. Bera,
\emph{Saturating quantum advantages in postselected metrology with the
positive Kirkwood--Dirac distribution},
Physical Review A \textbf{107} (2023), 042413,
\href{https://doi.org/10.1103/PhysRevA.107.042413}{doi:10.1103/PhysRevA.107.042413}.
\bibitem{lca}
M. Spriet, \emph{Characterizing the Kirkwood--Dirac positivity on second
countable LCA groups},
\href{https://arxiv.org/abs/2507.23628v3}{arXiv:2507.23628v3}, p.~7.
\bibitem{generic}
C. Langrenez, W. Salmon, S. De Bi\`evre, J. J. Thio, C. K. Long and
D. R. M. Arvidsson-Shukur,
\emph{Almost no experiments have classical Kirkwood--Dirac representations},
\href{https://arxiv.org/abs/2405.17557v2}{arXiv:2405.17557v2}
(2 July 2026), Theorem 1, (7).
\end{thebibliography}
\endgroup
\medskip\noindent\small
License: Creative Commons Attribution 4.0 International (CC BY 4.0).
\end{document}
